In their apparent motions, the Moon and the other large planets in our solar
system do not stray far from the plane of the ecliptic. Their apparent
trajectories described on the celestial sphere remain contained in a region
that extends symmetrically on both sides of the ecliptic, having a total width
of about $18^\circ$. Figure 1 shows the heliocentric orbit of the Earth in
the ecliptic plane, and Figure 2 the equivalent apparent geocentric orbit of
the Sun relative to the Earth.
\ioaafig{2022_TH_Q11-f1.pdf}
\ioaafig{2022_TH_Q11-f2.pdf}
The time intervals of the evolution of the true Sun in each of the 12
constellations of the Zodiac are as follows.
\begin{center}
\begin{tabular}{lll}
1) Aries & 21 III--20 IV & 2) Taurus, 21 IV--20 V\\
3) Gemini & 21 V--20 VI & 4) Cancer, 21 VI--22 VII\\
5) Leo & 23 VII--22 VIII & 6) Virgo, 23 VIII--22 IX\\
7) Libra & 23 IX--22 X & 8) Scorpio, 23 X--21 XI\\
9) Sagittarius & 22 XI--21 XII & 10) Capricorn, 22 XII--21 I\\
11) Aquarius & 22 I--19 II & 12) Pisces, 20 II--21 III
\end{tabular}
\end{center}
% sic: The source calls these intervals constellations although the dates are tropical zodiac sign intervals.
The International Olympiad on Astronomy and Astrophysics for Juniors,
Edition I, was to take place in Suceava, in Romania, on March 28--April 3,
2020. The event, proposed by Romania, was established at the International
Olympiad of Astronomy and Astrophysics, 13th Edition, held in Hungary on
August 2--10, 2019. But, given the known conditions, IOAA-jr, Edition I, is
taking place in Romania only now, in November 2022!

Hint: All tasks requiring direct measurements on drawings will be performed
on the enlarged pictures in source pages 5 (Figure 1) and 6 (Figure 2),
reproduced above. Both pages are part of the answer sheets.
\begin{description}
\item[(A)] Identify the following. \hfill[30 points]
  \begin{enumerate}
  \item The sign of the OIAA development, 13th Edition, Hungary, August 2--10, 2019.
  \item The sign when the OIAA-J was to take place, Edition I, March 28--April 3,
    2020, Romania (Edition I, proposal).
  \item The sign of the OIAA-J, Edition I, November 2022.
  \end{enumerate}
\item[(B)] Each specified Olympic event can be considered to have taken place
  on the day when the Sun, viewed from Earth, was in the middle of the angular
  range corresponding to its sign. \hfill[20 points]
  \begin{enumerate}
  \item Determine the time interval between any two of the events,
    $\Delta t_{H-R,1}$, $\Delta t_{R,1-R,2}$, and $\Delta t_{H-R,2}$.
  \item Estimate by direct measurements with a protractor the angles between
    the Earth's position vectors $\vec r_H$, $\vec r_{R,1}$, $\vec r_{R,2}$
    relative to the Sun and the apse line Aph--Ph, namely $\alpha_H$,
    $\alpha_{R,1}$, and $\alpha_{R,2}$.
  \end{enumerate}
\item[(C)] Determine the area swept out by the Earth's position vector relative
  to the Sun as follows. \hfill[20 points]
  \begin{enumerate}
  \item From $\vec r_{\mathrm{Hungary}}$ to $\vec r_{\mathrm{Romania},1}$.
  \item From $\vec r_{\mathrm{Romania},1}$ to $\vec r_{\mathrm{Romania},2}$.
  \item From $\vec r_{\mathrm{Hungary}}$ to $\vec r_{\mathrm{Romania},2}$.
  \end{enumerate}
\item[(D)] Determine the following. \hfill[30 points]
  \begin{enumerate}
  \item The Earth--Sun distance $r_{R,1}$ on the day when IOAA for Juniors,
    Edition I, should have taken place in Suceava in 2020. The Earth's orbit
    has semimajor axis $a=149597500\,\mathrm{km}$ and semiminor axis
    $b=149580670\,\mathrm{km}$.
  \item The Earth's acceleration $a_{E,1}$ on that day, compared with the
    acceleration $g_{S,1}$ in the Sun's gravitational field at distance $r_{R,1}$.
  \item The components of the Earth's velocity $\vec v_{R,1}$ parallel and
    perpendicular to the ellipse's major axis, $v_\parallel$ and $v_\perp$.
  \end{enumerate}
\end{description}
The Earth's angular momentum relative to the Sun and the total mechanical
energy of the Earth--Sun system are
\[
L=M_Eb\sqrt{\frac{GM_S}{a}},\qquad E=-G\frac{M_EM_S}{2a}.
\]
It is given: $G=6.67\times10^{-11}\,\mathrm{N\,m^2\,kg^{-2}}$,
$M_S=1.989\times10^{30}\,\mathrm{kg}$, and the Earth's revolution period
$T=365.256$ days.
